\documentclass[12pt,notitlepage]{report}
% The shared notes preamble lives outside this repo.
\input{/Users/aryankulkarni/Documents/y3-notes/preamble.tex}

\title{Week 6: Poisson regression for goals}
\author{Aryan Kulkarni}
\date{MTH6134 group project}

\begin{document}
\maketitle

\paragraph{Findings.}
Ashford score twice as many goals at home as away, and 5.6\% more for each
place lower the opponent finished. The Poisson model fits. The strength of
both effects rests partly on one match, an 8-0 home win over the bottom club.
Previous injuries are not explained by any variable and are overdispersed.

\section{The model}

Goals in match $i$ are $Y_i \dist \text{Poisson}(\mu_i)$ with
\[
  \log \mu_i = \beta_0 + \beta_1\,\text{home}_i + \beta_2\,\text{rank}_i ,
\]
so $e^{\beta}$ is a rate ratio: the factor by which the expected number of
goals is multiplied. The model assumes independent matches and a variance
equal to the mean.

\begin{minted}{r}
source("../common.R")

fit_goals <- glm(goals_scored ~ home_away + opponent_rank,
                 family = poisson, data = matches)
out(summary(fit_goals), "summary-goals")
out(round(exp(cbind(rate_ratio = coef(fit_goals), confint(fit_goals))), 3),
    "rate-ratios")
\end{minted}

\inputminted{text}{output/summary-goals.txt}
\inputminted{text}{output/rate-ratios.txt}

The home rate ratio is exactly 2.00 (95\% interval 1.18 to 3.49): 40 goals in
23 home matches against 20 in 23 away. Because every opponent is met once at
home and once away, rank is balanced across venues and does not disturb this
ratio. Each place down the table multiplies expected goals by 1.056 (1.017 to
1.098). The fitted mean runs from 0.43 goals away to the champions up to 3.0
at home to the bottom club.

\begin{minted}{r}
grid <- expand_grid(home_away = levels(matches$home_away), opponent_rank = 1:24)
grid$goals_scored <- predict(fit_goals, grid, type = "response")

fig <- ggplot(matches, aes(x = opponent_rank, y = goals_scored,
                           colour = home_away)) +
  geom_point(size = 1.6, alpha = 0.8) +
  geom_line(data = grid) +
  scale_colour_manual(values = c(Away = orange, Home = blue), name = NULL) +
  labs(x = "Opponent's final rank", y = "Goals scored")
save_fig(fig, "fitted", height = 2.8)
\end{minted}

\begin{figure}[H]
  \centering
  \incfig{fitted}
  \caption{Goals scored with the fitted Poisson means. The point at 8 goals
    is the home match against the 24th-placed club.}
  \label{fig:fitted}
\end{figure}

\section{Does the model fit?}

\begin{minted}{r}
fit_inter <- update(fit_goals, . ~ . + home_away:opponent_rank)
set.seed(1)
simulated <- replicate(2000, {
  y <- rpois(nrow(matches), fitted(fit_goals))
  deviance(glm(y ~ home_away + opponent_rank, family = poisson, data = matches))
})
out(c(
  deviance = deviance(fit_goals),
  df = df.residual(fit_goals),
  p_chisq = pchisq(deviance(fit_goals), df.residual(fit_goals),
                   lower.tail = FALSE),
  p_simulated = mean(simulated >= deviance(fit_goals)),
  dispersion = sum(resid(fit_goals, "pearson")^2) / df.residual(fit_goals),
  p_interaction = anova(fit_goals, fit_inter, test = "Chisq")$`Pr(>Chi)`[2]
), "fit-check")
\end{minted}

\inputminted{text}{output/fit-check.txt}

The residual deviance is 49.8 on 43 degrees of freedom. The $\chi^2$
approximation to its distribution is poor when the means are this small, so
the deviance is also compared with 2000 data sets simulated from the fitted
model: neither $p$ value (0.22 and 0.44) signals lack of fit. The Pearson
dispersion estimate is 0.99, so the variance matches the mean. The home
advantage does not vary with opponent strength ($p = 0.55$ for the
interaction).

\begin{minted}{r}
q <- qresid(fit_goals)
cook <- cooks.distance(fit_goals)
cook_cut <- 8 / (nobs(fit_goals) - 2 * length(coef(fit_goals)))

base_fig("checks", mfrow = c(1, 3), height = 2.3)
plot(fitted(fit_goals), q, col = pt, cex = 1.3,
     xlab = "Fitted mean", ylab = "Quantile residuals")
abline(h = 0, lwd = 2)
qqnorm(q, main = "", col = pt, cex = 1.3)
qqline(q, lwd = 2)
plot(cook, type = "h", col = blue, xlab = "Match", ylab = "Cook's distance")
abline(h = cook_cut, lty = 2)
text(8, cook[8], labels = 8, pos = 4, cex = 0.8)
dev.off()

fit_drop <- update(fit_goals, data = matches[-8, ])
out(c(cook_threshold = cook_cut, cook_match_8 = cook[[8]],
      fitted_match_8 = fitted(fit_goals)[[8]],
      prob_8_or_more = ppois(7, fitted(fit_goals)[[8]], lower.tail = FALSE)),
    "match-8")
out(round(cbind(coef(summary(fit_drop))[, c(1, 4)],
                exp(cbind(rate_ratio = coef(fit_drop), confint(fit_drop)))), 3),
    "without-match-8")
\end{minted}

\begin{figure}[H]
  \centering
  \incfig{checks}
  \caption{Quantile residuals against fitted means, their QQ plot, and Cook's
    distances with the threshold $8/(n - 2p) = 0.2$.}
  \label{fig:checks}
\end{figure}

\inputminted{text}{output/match-8.txt}
\inputminted{text}{output/without-match-8.txt}

The residuals show no trend and follow the normal line apart from one point.
That point is match 8, the 8-0 win, with a Cook's distance of 0.64 against a
threshold of 0.2. The model expects 3.0 goals in that match and gives 8 or
more a probability of 0.012, which is unusual for one match and unremarkable
as the most extreme of 46.

Without match 8 the home rate ratio falls from 2.00 to 1.71 and the rank
ratio from 1.056 to 1.034, and neither is significant at 5\% ($p = 0.061$
and $0.105$). The match is a real result and stays in the data. The honest
statement is that the direction of both effects is stable and their
significance depends on one match.

\section{Goals conceded}

\begin{minted}{r}
fit_conceded <- glm(goals_conceded ~ home_away + opponent_rank,
                    family = poisson, data = matches)
out(round(cbind(coef(summary(fit_conceded))[, c(1, 2, 4)],
                exp(cbind(rate_ratio = coef(fit_conceded),
                          confint(fit_conceded)))), 3), "conceded")
out(c(dispersion = sum(resid(fit_conceded, "pearson")^2) /
        df.residual(fit_conceded),
      largest_cook = max(cooks.distance(fit_conceded))), "conceded-check")
\end{minted}

\inputminted[fontsize=\footnotesize]{text}{output/conceded.txt}
\inputminted{text}{output/conceded-check.txt}

The mirror image holds. Ashford concede about half as many goals at home
(rate ratio 0.52, interval 0.26 to 1.00) and 5\% fewer per place of opponent
rank (0.95, interval 0.91 to 0.99). No match is influential. The dispersion
estimate of 0.57 is below 1, so the counts vary less than a Poisson variable
would, which makes these standard errors conservative. Week 10 returns to
this.

\section{Previous injuries}

\begin{minted}{r}
fit_prev <- glm(prev_injuries ~ position + age + training_load + bmi,
                family = poisson, data = players)
explained <- fit_prev$null.deviance - deviance(fit_prev)
out(c(deviance_explained = explained, df = 6,
      p = pchisq(explained, 6, lower.tail = FALSE),
      dispersion = sum(resid(fit_prev, "pearson")^2) / df.residual(fit_prev)),
    "prev")
\end{minted}

\inputminted{text}{output/prev.txt}

The six coefficients together reduce the deviance by 4.8, no more than chance
($p = 0.58$): there is no evidence that anything we measure predicts how many
injuries a player had in the previous three seasons. The dispersion estimate
is 1.31, so the counts are more variable than Poisson, as the plot in week 3
suggested. Overdispersion makes a Poisson test too ready to reject, so the
absence of an effect is not an artefact of the wrong variance. A Poisson
model is the wrong description of this variable, and week 10 looks for a
better one.

\section{Summary}

\begin{itemize}
  \item \mintinline{r}|goals_scored ~ home_away + opponent_rank| with a log
    link fits well: home rate ratio 2.00, rank ratio 1.056 per place.
  \item One match carries much of the evidence, and the write-up should say
    so.
  \item Goals conceded are underdispersed and previous injuries
    overdispersed. Both go to week 10.
\end{itemize}

\end{document}
